% ---------------------------------------------------------------------
% Section 1: Introduction.
% Drafted by Claude following the paragraph plan in Outline.docx.
% Literature items the outline asked Claude to research are flagged
% with \Claude{} comments; candidates with summaries are in
% notes/literature-candidates.md for the authors to select from.
% ---------------------------------------------------------------------
\section{Introduction}\label{sec:intro}

Partially observed Markov decision processes (POMDPs) are pervasive in
operations research \citep{krishnamurthy2016partially}. In a POMDP, a
decision maker (DM) repeatedly chooses actions under uncertainty about
an underlying state, observing signals that partially reveal the state
and updating her beliefs over time. Examples include technology
adoption, where a consumer or firm decides whether to adopt a new
technology of uncertain benefit or wait and gather more information
\citep{mccardle1985information, lippman1987does, ulu2009uncertainty};
machine maintenance, where a DM decides when to inspect or replace
equipment whose condition is only partially observed
\citep{white1979optimal, maillart2006maintenance}; inventory control
with unobserved demand states \citep{treharne2002adaptive}; and
healthcare applications such as prostate biopsy referral decisions
\citep{zhang2012optimization} and patient-type adaptive treatment
planning \citep{skandari2021patient}. See
\citet{krishnamurthy2016partially} for many other examples.

Many POMDP models of interest have \emph{monotone} structure: the
optimal action increases as the DM's beliefs increase in a stochastic
sense. In technology adoption models, as beliefs about the benefits of
the technology improve, optimal policies move toward adoption
\citep{mccardle1985information, ulu2009uncertainty}. In machine
maintenance models, as beliefs about the state of the machine worsen,
optimal policies move toward replacement \citep{white1979optimal,
lovejoy1987some}. And as beliefs point toward cancer, optimal policies
move toward a biopsy \citep{zhang2012optimization}. A substantial
literature provides sufficient conditions for such structure:
\citet{lovejoy1987some} gives conditions for the optimal value
function of a finite POMDP to be monotone in the prior, with priors
ordered by likelihood ratios, for the optimal policy in a machine
replacement problem to be monotone, and, in the general case, for the
optimal policy to be bounded below by a monotone function.
\Claude{Your outline asks for additional structural-results
references, especially by Vikram Krishnamurthy. A candidate list with
summaries is in notes/literature-candidates.md for you to select
from; once you choose, I will fold them into this paragraph.}

How should one compare information sources in a POMDP? The classical
tool is Blackwell's comparison of experiments
\citep{blackwell1951comparison}: if, for each action, the signal
process is garbled by adding noise that is independent of the state,
the resulting POMDP has a lower optimal value than the original.
Establishing this value comparison uses the convexity of the POMDP
value function in the prior \citep{smallwood1973optimal} and requires
no monotonicity of the value function or optimal policies. Blackwell's
order is powerful precisely because it is preference free---a garbling
is worse in \emph{every} decision problem---but this universality
makes it a very strong requirement, and few pairs of information
structures can be ranked by it \citep{lehmann1988comparing,
kim2023comparing}. A large literature has therefore sought weaker
orders that rank more information structures by restricting the class
of decision problems, most prominently to \emph{monotone} decision
problems: \citet{lehmann1988comparing} introduces the accuracy order
for single-decision problems with monotone likelihood-ratio signals;
\citet{persico2000information} and \citet{quah2009comparative} extend
the reach of Lehmann's comparison to broader classes of monotone
payoffs; and \citet{jewitt2007information} connects Lehmann's order to
Blackwell dominance on dichotomies. \citet{athey2018value} study
monotone decision problems in a related but different setup, where
information structures are compared for a \emph{fixed} prior belief:
their orderings are prior dependent, which permits finer comparisons
when the prior is known, whereas the orders we study---like Blackwell's
and Lehmann's---rank signal processes uniformly across priors.
\Claude{Two items to verify here: (i) the outline asks for citations
supporting the ``Blackwell is too strong'' criticism---I cite Lehmann
and Kim, both of whom make the point explicitly; add others if you
have favorites. (ii) The outline asks whether Krishnamurthy has papers
using Lehmann's order: I could not identify any---his structural
POMDP results use Blackwell dominance of observation kernels (see
notes/literature-candidates.md)---please confirm before we assert
anything in text.}

All of these weaker orders, however, were developed for static,
single-decision problems. In this paper, we ask the corresponding
question for dynamic problems: when is one signal process worth more
than another in a \emph{monotone POMDP}? Because we focus on a smaller
class of POMDPs, we can obtain a weaker order among information
structures that is implied by Blackwell sufficiency. We say a signal
process $\iX$ is \emph{LR-better} than a signal process $\iY$ if $\iY$
can be obtained from $\iX$ by adding a state-dependent garbling such
that the garbling distribution at a lower state likelihood-ratio (LR)
dominates the garbling distribution at a higher state---that is, the
added noise systematically pushes signals upward in low states and
downward in high states, degrading the signal's ability to
discriminate between states. Our order is a likelihood-ratio
strengthening of the \emph{monotone quasi-garbling} order of
\citet{kim2023comparing}, who studies monotone decision problems with
a single decision and shows that monotone quasi-garbling---where the
garbling distribution at a lower state first-order stochastically
dominates the garbling at a higher state---is a necessary and
sufficient condition for all decision makers with monotone decision
problems to obtain higher expected payoffs. LR dominance implies
first-order stochastic dominance, so the LR-better order implies
monotone quasi-garbling. The strengthening is exactly what the dynamic
setting demands: LR dominance survives Bayesian updating and state
transitions in a POMDP, whereas first-order stochastic dominance does
not \citep[see][for a discussion]{ulu2009uncertainty}.

Our main result (Theorem~\ref{thm:main}) shows that in a monotone
POMDP---one with monotone rewards, monotone likelihood-ratio signal
and transition processes, and supermodular structure delivering
monotone optimal policies---an LR-better signal process yields a
higher optimal value function at every horizon. We also map the
relationships among the orders (Section~\ref{sec:orders}): under the
monotone likelihood ratio property, Blackwell's order implies the
LR-better order, which implies monotone quasi-garbling, which is in
turn equivalent to Lehmann's accuracy order. We illustrate the
framework with technology adoption, machine replacement, and medical
decision-making models, present Normal, Poisson, and Uniform families
of examples in which LR-better comparisons arise naturally, and show
by example that when the state dependence of the noise runs in the
opposite direction---signals from better states are \emph{more} likely
to be observed---the two signal processes may be genuinely
incomparable, with value functions that cross as a function of the
prior.

The rest of the paper is organized as follows.
Section~\ref{sec:model} presents the monotone POMDP model and its
structural properties. Section~\ref{sec:comparison} introduces the
LR-better order and establishes the main value-comparison theorem.
Section~\ref{sec:orders} relates the LR-better order to the Blackwell,
Lehmann, and monotone quasi-garbling orders.
Section~\ref{sec:applications} presents applications and examples,
and Section~\ref{sec:conclusion} concludes.
